\chapter{Magic and Special Abilities}
łabel{ch:magic}

Magic in this game is powerful but dangerous−−−a negotiation with reality itself that always carries risks. This chapter covers the core magical systems and the distinct paths a character may follow.

\begin{tcolorbox}[colback=black!3!white,colframe=black!40!white,title={Sidebar: \texttt{[TAGS]} & Magic},enhanced]
∈dex{TAGS}∈dex{magic!TAGS}
  \textbf{What are \texttt{[TAGS]}?} Effects in \textit{Fate's Edge} are communicated via \texttt{[TAGS]}. Each \texttt{[TAG]} is a discrete effect gated behind a Talent, Rite, spell, or asset−−−\emph{it cannot be invoked spontaneously} unless a rule grants access.
  
  \medskip
  \textbf{How they're used.} \texttt{[TAGS]} provide a common language for describing effects, especially when players invent spells via \emph{Free Casting}. Many prewritten spells and abilities also list their \texttt{[TAGS]} for clarity.
\end{tcolorbox}

\section{The Nature of Magic}
łabel{sec:magic−nature}
∈dex{magic!nature}

Magic is not a safe tool but a dangerous force:
\begin{itemize}
ℹtem \textbf{Powerful}: Can reshape battles, stories, or even the world.
ℹtem \textbf{Controlled}: Every use generates \textbf{Story Beats (SB)} that manifest as backlash (see \ref{ch:core−mechanics}).
ℹtem \textbf{Thematic}: Effects and consequences align with the type of magic used.
ℹtem \textbf{Volatile}: Never fully predictable or controllable.
ℹtem \textbf{Narrative}: Casting is always a significant story moment.
\end{itemize}

\paragraph{Table Vignette:}
\emph{``I can hold the avalanche,'' says Mira, fingers trembling. ``But something will answer.''} The party nods−−−risk accepted, stakes clear.

\subsection{Patron as Co−Star}
łabel{sec:patron−co−star}
∈dex{Patrons!active vs passive}

Your Patron is not just a source of power−−−they are a \textbf{co−star} in your story. How often they appear depends on your campaign's length. Ask your GM which style they will use:

\begin{description}
  ℹtem[Passive Patrons (long campaigns)] −− The Patron stays in the background, appearing only when Obligation thresholds are breached. Their demands feel like slow, inevitable consequences.
  ℹtem[Active Patrons (short campaigns / one‑shots)] −− The Patron appears often: in dreams, omens, or through intermediaries. They offer bargains, push the Obligation spiral, and make the price of power immediate.
\end{description}

\noindent\textbf{Player advice:} Embrace Patron demands as \emph{story hooks}, not punishments. When the GM says ``Your Patron requires a task,'' lean into it. That task will lead to adventure, not just a tax. A Patron who intervenes is a Patron who cares −− and a caring Patron is a source of power and drama.

\begin{tcolorbox}[colback=gray!5!white, colframe=gray!60!white, title=For Full Patron Details]
∈dex{Patrons!full list}
The complete descriptions of all Patrons−−−including their Rites, Obligation mechanics, and corruption tables−−−are found in \textbf{Chapter \ref{ch:world−powers−obligation} (World Powers and Obligation)}. That chapter also covers terrestrial patrons, faction obligations, and the mechanics of clearing debt. This chapter focuses on the \emph{paths} of magic; for the \emph{sources} of power, see the World Powers chapter.
\end{tcolorbox}

\section{Paths of Magic}
łabel{sec:paths−of−magic}
∈dex{magic!paths}

Magic in \textit{Fate's Edge} is not a single system. Characters may follow one of several distinct paths, each with its own mechanics, costs, and narrative flavor. The five primary paths are:

\begin{description}
  ℹtem[Free Caster] Improvisational magic using \texttt{[TAGS]}. Flexible but risky (\ref{sec:free−caster}).
  ℹtem[Runekeeper] Structured rites granted by a Patron. Reliable but accrues Obligation (\ref{sec:runekeeper}).
  ℹtem[Invoker] Ritual magic using Patron symbols. Safe but slow (\ref{sec:invoker}).
  ℹtem[Cantor] Songs that mimic Low Rites. Accessible but corrupting (\ref{sec:cantor}).
  ℹtem[Summoner] Binding and commanding spirits. Powerful but requires control (\ref{sec:summoner}).
\end{description}

A character may dabble in multiple paths, but specialization is more effective.

\subsection{Free Caster (Spellcraft)}
łabel{sec:free−caster}
∈dex{magic!free casting}∈dex{Spellcraft}

Free casters shape raw magic through will and knowledge. They do not rely on Patrons or rote rites.

⫕section*{Requirements}
\begin{itemize}
  ℹtem Talent: \textbf{Spellcraft} (6 XP) −− see \ref{ch:talents}.
  ℹtem Appropriate magical skill (Arcana, Elemental Magic, etc.).
\end{itemize}

⫕section*{Casting Procedure}
łabel{sec:free−casting−procedure}
\begin{enumerate}
  ℹtem \textbf{Declare Intent:} Describe the desired effect using \texttt{[TAGS]}.
  ℹtem \textbf{Choose Approach:} Attribute + magical skill.
  ℹtem \textbf{Set Position & DV:} DV = 1 + number of \texttt{[TAGS]}. Dangerous tags add +2 each.
  ℹtem \textbf{Roll:} Resolve normally (see \ref{ch:core−mechanics}). Each 1 generates a Story Beat.
\end{enumerate}

For powerful or complex effects, use the \textbf{Casting Loop} (two actions: Weave then Cast). Simple effects (one tag) may be cast as a \textbf{cantrip} (one action, no separate Weave).

⫕section*{Backlash}
łabel{sec:free−caster−backlash}
On a failure or mixed result, the GM may impose:
\begin{itemize}
  ℹtem Harm 2 (Arcane)
  ℹtem +2 Fatigue
  ℹtem Corruption +1
  ℹtem A catastrophic side effect
\end{itemize}
Backlash is always flavored by the elements used.

⫕section*{Free Caster Talents}
łabel{sec:free−caster−talents}
∈dex{Talents!free caster}
\begin{description}
  ℹtem[Ritual Mastery (6 XP)] Reduce casting time for major spells by half.
  ℹtem[Elemental Focus (4 XP)] Gain +1 die when using a chosen element.
  ℹtem[Stabilize Weave (8 XP)] Once per session, cancel 1 SB from a failed cast.
\end{description}

\subsection{Runekeeper (Rites)}
łabel{sec:runekeeper}
∈dex{magic!runekeeper}∈dex{Rites}∈dex{Obligation}

Runekeepers bind themselves to a single Patron, learning structured Rites that produce reliable effects.

⫕section*{Requirements}
\begin{itemize}
  ℹtem \textbf{Thiasos (Familiar)} (2 XP) −− a token or bond with the Patron (see \ref{ch:talents}).
  ℹtem \textbf{Codex} (4 XP) −− a personal record of Rites.
  ℹtem \textbf{Patron Choice:} Select one Patron (e.g., Traveler, Oath of Flame, Morag).
\end{itemize}

⫕section*{Invoking a Rite}
łabel{sec:invoking−rite}
\begin{itemize}
  ℹtem \textbf{Action:} 1 action.
  ℹtem \textbf{Cost:} Mark the Rite's listed Obligation (e.g., \texttt{Obligation: 6}). This cost is reduced by 1 if you display your Patron's Symbol (Invoker only).
  ℹtem \textbf{Effect:} The Rite works as written.
  ℹtem \textbf{DV for Rites that require a roll:} 
  \[
  \text{DV} = \max(\text{Obligation Cost} − \text{Spirit},  \text{Rite Tier})
  \]
  where Rite Tier is 1 for Low, 2 for Standard, 3 for High/Epic.
  ℹtem \textbf{Push It:} Once per scene, mark +1 additional Obligation to amplify the effect (+1 die or +1 Effect).
\end{itemize}

⫕section*{Patron's Gift (Imbuement)}
łabel{sec:patrons−gift}
Once per scene, as 1 action, imbue a held item with +1 Melee and +1 Thematic skill (e.g., Navigation for Traveler). Lasts one scene. Mark +1 Obligation.

⫕section*{Runekeeper Talents}
łabel{sec:runekeeper−talents}
∈dex{Talents!runekeeper}
\begin{description}
  ℹtem[Efficient Invocation (4 XP)] Reduce Obligation cost of one specific Low Rite by 1 (min 0).
  ℹtem[Patron's Favor (8 XP)] Once per session, ignore Obligation cost for a Rite.
  ℹtem[Runebound Resilience (18 XP)] Convert up to 2 Harm into Obligation instead of damage.
\end{description}

\subsection{Invoker (Symbols)}
łabel{sec:invoker}
∈dex{magic!invoker}∈dex{Symbols}

Invokers do not bond permanently. They use consecrated Symbols to borrow a Patron's power through ritual.

⫕section*{Requirements}
\begin{itemize}
  ℹtem \textbf{Patron's Symbol} (4 XP per Patron). One Symbol per Patron.
  ℹtem No Thiasos or Codex required.
\end{itemize}

⫕section*{Ritual Invocation}
łabel{sec:ritual−invocation}
\begin{itemize}
  ℹtem \textbf{Time:} Significant Time (DV in rounds). +1 round per DV.
  ℹtem \textbf{Cost:} Mark the Rite's listed Obligation.
  ℹtem \textbf{Symbol Benefit:} If the Symbol is openly displayed, reduce DV by 1 (min 1) and reduce Obligation cost by 1 (min 0).
  ℹtem \textbf{No Symbol:} +1 DV, +1 Obligation, and on a partial/failure, GM gains +1 extra SB.
\end{itemize}

⫕section*{Crack the Seal}
łabel{sec:crack−the−seal}
Resolve a known Rite instantly as 1 action. Set the Symbol to \textbf{Compromised} and mark double the Rite's base cost (minimum +2, or +3 for High Rites). Restore the Symbol during downtime (Arcana test DV 3) or spend 1 XP.

⫕section*{Invoker Talents}
łabel{sec:invoker−talents}
∈dex{Talents!invoker}
\begin{description}
  ℹtem[Borrowed Grace (Free)] Spend 1 Boon or 1 Fatigue to gain +1 Melee or +1 Thematic for one action. Mark +1 Obligation.
  ℹtem[Crack Specialist (6 XP)] Reduce Obligation cost of Crack the Seal by 1.
  ℹtem[Symbol Resonance (4 XP)] When using a Symbol, gain +1 die to the ritual roll.
\end{description}

\subsection{Cantor (Songs)}
łabel{sec:cantor}
∈dex{magic!cantor}∈dex{Songs}∈dex{Corruption}

Cantors perform Low Rites as songs, using voice and performance instead of Patron bonds.

⫕section*{Requirements}
\begin{itemize}
  ℹtem Talent: \textbf{Cantor's Path} (8 XP) −− see \ref{ch:talents}.
  ℹtem Prerequisites: Lore 1+, Performance 2+, Presence 2+.
\end{itemize}

⫕section*{Performing a Song}
łabel{sec:performing−song}
\begin{itemize}
  ℹtem \textbf{Action:} 1 action to begin; resolves at start of next turn.
  ℹtem \textbf{Test:} Lore + Performance (DV 2−−3).
  ℹtem \textbf{Cost:} Materials as listed. No Obligation on success.
  ℹtem \textbf{Corruption:} Track a Corruption Timer (size = Spirit). Mark segments when you Push, perform a Resonant Rite, or when the GM spends SB on your occult activities.
\end{itemize}

⫕section*{Push It}
łabel{sec:cantor−push}
Resolve immediately. Mark Fatigue 1, mark toward Corruption accumulation, and the GM gains 1 SB (social/Patron fallout).

⫕section*{Corruption Timer }
łabel{sec:corruption−timer}
When the timer fills, gain a thematic benefit and drawback from the last Patron whose Rite you performed. Then \textbf{set your current Corruption segments to your Tier} (minimum 1). The timer's maximum size remains equal to your Spirit. This means you are never completely free of Corruption; high‑Tier Cantors always carry a remnant of past blooms.

⫕section*{Resonant Rites (Founding a Chorus)}
łabel{sec:resonant−rite}
A \textbf{Resonant Rite} is any Cantor song performed with dramatic intent, witnessed by at least 10 people, and dedicated to a specific Patron. The GM determines if a performance qualifies.

⫕section*{Chorus Cults (Optional)}
łabel{sec:chorus−cults}
If you perform a Resonant Rite before 10+ witnesses, you may found a \textbf{Chorus}. This grants a minor follower benefit (see \ref{ch:assets−followers}), but you must mark Obligation to the cult's patron (the patron whose Resonant Rite was sung to form the cult) each session you use it. \textbf{If neglected, the cult may fracture} −− your followers could splinter into a rival sect, spread false rumours, or ally with a Patron's enemy. Keep an eye on your followers; a neglected Chorus can become a serious problem.

\textbf{Example minor follower benefits (work with your GM):}
\begin{itemize}
  ℹtem Once per session, gain +1 die on a Performance roll or on a social roll where the crowd's opinion matters.
  ℹtem Once per session, call upon the cult to provide a single piece of information from within a community where the cult has a presence.
  ℹtem Once per session, gain +1 Boon when you perform a Song in a public space.
\end{itemize}

⫕section*{High Cantor and Instant Songs}
łabel{sec:high−cantor}
The \textbf{High Cantor} talent (below) allows you to learn Standard Rites as songs and resolve them instantly (no delay). When you sing a Standard Rite as a Song, you may also \textbf{Push} it for additional effect (e.g., +1 Effect) at the usual cost of 1 Fatigue and marking a Corruption segment.

⫕section*{Cantor Talents}
łabel{sec:cantor−talents}
∈dex{Talents!cantor}

\begin{description}
  ℹtem[Resonant Performance (3 XP)] Gain +1 die when performing for an audience of 5+.

  ℹtem[Song Weaver (4 XP)] Combine two compatible Songs for +1 Effect. Two Songs are compatible if they share a thematic tag (e.g., both affect healing, positioning, spirits) or if they are from the same Patron. The GM has final say. Sing both as a single action; resolve the Performance roll once; if successful, both effects occur with +1 Effect.

  ℹtem[High Cantor (18 XP)] Learn Standard Rites as songs; resolve instantly (as described above).

  ℹtem[Steady Voice (2 XP)] \textbf{Passive.} When you Push a Song, you may choose to mark 1 additional Fatigue instead of marking a Corruption segment (once per scene). You cannot use this talent if your Corruption timer is already full. This talent gives low‑Spirit Cantors an emergency buffer at the cost of greater exhaustion. \textit{Requires: Cantor's Path.}

  ℹtem[Patron's Dedication (4 XP)] \textbf{Prerequisite:} Cantor's Path. \\
  \textbf{Effect:} You bind your voice to a single Patron (choose one when you take this talent). Your Corruption timer is always keyed to that Patron −− it never resets to a different Patron's table when you bloom. Whenever you would gain Corruption from any source, you mark it on this same timer. \\
  \textbf{Downside (Rivalry):} Singing a Song (Low or Standard Rite) of any other Patron is treacherous. When you do so, your Position for that song is worsened by one step (Dominant \(→\) Controlled, Controlled \(→\) Desperate). This penalty applies even if the song would otherwise be instant or free.
\end{description}

\subsection{Summoner (Pact−Whisperer)}
łabel{sec:summoner}
∈dex{magic!summoning}∈dex{Leash}∈dex{Pact−Whisperer}

Summoners call and bind Outsiders for temporary aid.

⫕section*{Requirements}
\begin{itemize}
  ℹtem Talent: \textbf{Pact−Whisperer} (2 XP) −− see \ref{ch:talents}.
  ℹtem Lesser Pactwright (2 XP) for Cap 1 spirits; Greater Pactwright (4 XP) for Cap 3.
\end{itemize}

⫕section*{Core Mechanics}
łabel{sec:summoning−mechanics}
\begin{itemize}
  ℹtem \textbf{Call (1 action):} Manifest a spirit at Near range.
  ℹtem \textbf{Bind:} Spend 1 Boon or mark 1 Fatigue.
  ℹtem \textbf{Leash Capacity:} Cap + Spirit segments.
  ℹtem \textbf{Tick Leash:} Spirit takes harm, commanded against nature, crosses a \texttt{[WARD]}, or you split focus.
  ℹtem \textbf{Departure:} When Leash fills, spirit acts once according to its nature, then departs (or turns hostile at GM discretion).
\end{itemize}

⫕section*{Spirit Whims Table (When the Spirit Acts on Its Own)}
łabel{sec:spirit−whims}
When the Leash fills or the spirit acts independently, the GM may choose (or roll) from this table to make the action interesting −− not just damage.

\begin{center}
\begin{tabular}{cl}
⊤rule
\textbf{d6} & \textbf{Spirit's Independent Action} \\
∣rule
1 & \textbf{Reveal a Secret:} The spirit blurts out a hidden fact about a PC or NPC −− true, embarrassing, or dangerous. \\
2 & \textbf{Break Something:} The spirit shatters a mundane object (a lantern, a rope bridge, a crucial tool). \\
3 & \textbf{Lead Astray:} The spirit moves a short distance and creates a false trail, forcing the party to waste time. \\
4 & \textbf{Call for Help:} The spirit summons a minor ally (a rat swarm, a will‑o'‑wisp, a lesser elemental). \\
5 & \textbf{Curse a PC:} The spirit leaves a lingering mark (−−1 die on one skill until the next sunrise). \\
6 & \textbf{Bargain:} The spirit offers a deal: ``Release me now, and I'll tell you where the treasure is hidden.'' \\
⊥tomrule
\end{tabular}
\end{center}

⫕section*{Boon Finesse}
łabel{sec:boon−finesse}
Once per round, spend 1 Boon to clear 1 Leash tick.

⫕section*{Summoner Talents}
łabel{sec:summoner−talents}
∈dex{Talents!summoner}
\begin{description}
  ℹtem[Spirit Synergy (4 XP)] When commanding two spirits, reduce each Leash by 1.
  ℹtem[Spirit Link (10 XP)] Spirits act on your turn; commands are free actions.
  ℹtem[True Name Keeper (15 XP)] Call any previously encountered spirit by true name; reduce Leash Capacity by 2.
\end{description}

\section{Universal Mechanics}
łabel{sec:universal−magic−mechanics}

\subsection{Obligation Capacity}
łabel{sec:obligation−capacity}
∈dex{Obligation!capacity}

A character's \textbf{Obligation Capacity} equals Spirit + Presence.  
Track total Obligation segments across all Patrons.

\begin{itemize}
  ℹtem \textbf{Exceeding Capacity:} For each segment above Capacity, mark 1 Fatigue. Cannot invoke Rites until Obligation is reduced.
  ℹtem \textbf{Resolution:} Reduce Obligation through downtime service, Patron tasks, or ritual cleansing.
\end{itemize}

\noindent\textbf{Player guidance:} You track your own Obligation (see Player‑Managed Modules, \ref{sec:player−managed−modules}). When a Patron intrudes (e.g., a demand or omen), treat it as a \textbf{plot hook, not a penalty}. The Patron is offering you a story.

\subsection{Difficulty Value for Rites}
łabel{sec:dv−rites}
When a Rite (Runekeeper or Invoker) calls for a roll, the GM sets DV using:
\[
\text{DV} = \max(\text{Obligation Cost} − \text{Spirit},  \text{Rite Tier})
\]
\noindent\textit{Rite Tier:} Low = 1, Standard = 2, High/Epic = 3.\\
Example: A Low Rite (Tier 1) with Obligation cost 4, Spirit 2 → DV = max(4−2, 1) = 2.

\begin{tcolorbox}[colback=green!5!white,colframe=green!75!black,title=Components as Devotion]
∈dex{rituals!components}∈dex{Obligation!reduction}
Even if you already have a Patron's Symbol, Codex, or Thiasos, voluntarily acquiring rare components for a ritual shows dedication. Use this rule of thumb:

\begin{itemize}
  ℹtem \textbf{Ritual with no components (just Symbol/Codex):} Normal DV, normal Obligation cost.
  ℹtem \textbf{Ritual with common components (herbs, salt, a candle):} Reduce Obligation by 1 (minimum 0).
  ℹtem \textbf{Ritual with rare or dangerous components (a phoenix feather, a drop of troll blood, a whispered secret):} Reduce DV by 1 \textbf{and} clear 1 segment of existing Obligation to that Patron (once per session max).
  ℹtem \textbf{Ritual with a truly legendary component (the crown of a dead king, a tear from a goddess):} The Patron may waive all Obligation for that ritual −− and even grant a free Boon.
\end{itemize}
This turns component hunting into adventure hooks and gives you a meaningful choice: use the easy, costly path or invest legwork to earn the Patron's favor.
\end{tcolorbox}

\subsection{Backlash}
łabel{sec:magic−backlash}
∈dex{Backlash}

Any magical roll can generate backlash. The GM spends Story Beats (SB) to introduce complications tied to the element or Patron involved.

\begin{center}
\small
\begin{tabular}{ll}
⊤rule
\textbf{Result} & \textbf{Backlash Severity} \\
∣rule
Partial & Minor (e.g., Position −1, Fatigue 1, minor complication) \\
Miss & Major (e.g., Harm 1, condition, significant complication) \\
Critical failure (many 1s) & Catastrophic (e.g., Patron intervention, scene shift) \\
⊥tomrule
\end{tabular}
\end{center}

\begin{tcolorbox}[title=Backlash Severity by Position, colback=white!97!gray, colframe=black!80!gray]
\begin{tabular}{ll}
⊤rule
\textbf{Position} & \textbf{Backlash Severity on Partial/Miss} \\
∣rule
Dominant & Minor −− one SB complication, manageable cost. \\
Controlled & Moderate −− 1−−2 SB complication, temporary setback. \\
Desperate & Severe −− 2−−3 SB complication, lasting consequences or Harm. \\
⊥tomrule
\end{tabular}
\end{tcolorbox}

⫕section*{Elemental Backlash Examples}
łabel{sec:elemental−backlash}
∈dex{Backlash!elemental}
\begin{itemize}
  ℹtem \textbf{Fire:} Burns, flares, smoke.
  ℹtem \textbf{Water:} Slick footing, gear corrosion, flooding.
  ℹtem \textbf{Earth:} Tremors, encumbrance, binding roots.
  ℹtem \textbf{Air:} Dispersal, deafening gusts, thrown debris.
  ℹtem \textbf{Fate:} Opportunities close, predictions become set.
  ℹtem \textbf{Life:} Uncontrolled growth, fever, exhaustion.
  ℹtem \textbf{Luck:} Fortunes reverse at the worst moment.
  ℹtem \textbf{Death/Dreams:} Whispers, numbness, portents.
\end{itemize}

\subsection{DV by Spell Scope}
łabel{sec:dv−spell−scope}
Use this table as a quick reference when the GM asks for a DV for a new spell:

\begin{center}
\begin{tabular}{lll}
⊤rule
\textbf{Scope} & \textbf{Example} & \textbf{DV} \\
∣rule
Cantrip & Light a candle, sense magic & 2 \\
Combat & Fire bolt, force shield & 3 \\
Ritual (small) & Find a lost object, speak with a spirit & 4 \\
Ritual (large) & Summon a storm, teleport across a region & 5+ \\
Epic & Raise a dead army, rewrite a law of reality & 6+ (with quest) \\
⊥tomrule
\end{tabular}
\end{center}

\begin{tcolorbox}[colback=yellow!5!white,colframe=yellow!75!black,title=When in Doubt, Measure Tags by Narrative Impact]
∈dex{TAGS!measuring by narrative impact}
A spell's power is not measured by how many dice it rolls, but by \textbf{how much it changes the story}. Use this rule of thumb:

\begin{itemize}
  ℹtem \textbf{Minor change (1 TAG):} The effect affects only yourself or a single object. Examples: \texttt{Light}, \texttt{[HEAL]}, \texttt{Veil} (self only). DV 2−−3.
  ℹtem \textbf{Moderate change (2 TAGS):} The spell affects a whole zone, creates a significant advantage, or bypasses a notable obstacle. Examples: \texttt{Area} + \texttt{Force}, \texttt{Barrier} across a doorway. DV 3−−4.
  ℹtem \textbf{Major change (3+ TAGS):} The spell rewrites the situation, solves a major problem outright, or introduces a new faction/ally/threat. Examples: \texttt{Teleport} across a city, \texttt{Summon} a Cap 3 spirit, \texttt{Transform} a person into a beast. DV 5+ and often a ritual.
\end{itemize}
If your description would \emph{skip over a scene} (``I teleport us past the entire fortress''), that's a major change −− require components, a timer, or a steep cost. If it would only \emph{add flavour} (``I make my eyes glow red''), that's free. Let the fiction guide the DV, not a pre‑computed formula.
\end{tcolorbox}

\subsection{Example: Accepting a Cost −− The Hard Choice}
łabel{sec:magic−choice−example}
When a magical roll generates SB, the GM may offer you a choice instead of a flat consequence.

\paragraph{Example:}
You cast a ritual to divine an enemy's location. The roll generates 2 SB. The GM says:
\begin{quote}
``The vision comes clearly −− but your Patron's symbol glows hot on your hand. \textbf{You can have the vision for free, or you can let the spell fizzle and keep your Obligation unchanged.} If you accept the vision, mark +2 Obligation and your Patron will demand a task next session. What do you choose?''
\end{quote}
This turns a success into a \emph{narrative fork}. Embrace it −− the Patron's task will be an adventure hook.

\section{Free Casting (TAGS System)}łabel{sec:free_castingₜags}
∈dex{Freeform magic!TAGS}∈dex{Tags}

Some casters shape raw forces through shared arcane grammar known as \textbf{TAGS}. Construct a spell at the table using a short phrase of TAGS.

\subsection{Spell Structure}
\textbf{Intent} + \textbf{Target} + \textbf{Tags} = effect. Example: \textit{``I unleash Burning • Area • Force against the marauders.''}

\subsection{Base Difficulty Value (DV)}
\[
\text{DV} = 1 + \text{number of TAGS}
\]
Dangerous TAGS (Teleportation, Transformation, Domination) add +2 DV instead of +1.

\subsection{Casting Roll}
Roll \textbf{Wits + Arcana} (or similar). Success = spell works. Failure or any 1 = Backlash.

\subsection{Backlash for Free Casters}∈dex{Backlash!Free Caster}
Choose one:
\begin{itemize}
    ℹtem Harm 2 (Arcane)
    ℹtem +2 Fatigue
    ℹtem Corruption +1
    ℹtem Catastrophic side effect (GM describes)
\end{itemize}
If the spell includes a ``Dangerous'' TAG, Backlash triggers on \emph{mixed} results as well.

\subsection{TAG Library (Abridged)}łabel{subsec:tagₗibrary}
∈dex{Tags!List}

\begin{longtable}{p{3.5cm} p{10cm}}
∩tion{Common TAGS}łabel{tab:tagsₐbridged}\\
⊤rule
\textbf{TAG Category} & \textbf{Examples} \\
∣rule
Elemental & Burning, Freezing, Storm, Stone, Wave, Wind \\
Force & Force, Area, Strike, Wall, Bind, Dispel \\
Mind & Veil & Veil, Scry, Memory, Command, Fear \\
Life & Body & \texttt{[HEAL]}, Purify, Strengthen, Waken, Beast \\
Space & Motion (dangerous, +2 DV) & Fold, Gate, Gravity \\
Creation & Transformation (dangerous, +2 DV) & Create, Summon, Transmute, Animate \\
⊥tomrule
\end{longtable}

\textit{Note:} \texttt{Leap} (short‑range teleport) is considered a dangerous TAG; it adds +2 DV. The SB cost should be less dramatic than [TELEPORT] or [GATE], though.

\subsection{Magic in Social Situations}
łabel{sec:magic−social}
∈dex{magic!social use}
Using magic socially carries additional risks:
\begin{itemize}
  ℹtem \textbf{Discretion:} Hidden casting may require a Stealth or Subterfuge test.
  ℹtem \textbf{Consent:} Affecting minds without permission is often considered a crime or taboo.
  ℹtem \textbf{Reputation:} Known magic users may face distrust or fear.
  ℹtem \textbf{Legal Consequences:} Many cities restrict magic use without a license.
\end{itemize}

\section{Magical Item Creation}
łabel{sec:magic−item−creation}
∈dex{magic!item creation}

Creating permanent magical items is a complex, campaign−level endeavor. See the full rules in the \textit{Game Master's Guide}. In brief:

\begin{itemize}
  ℹtem Requires high magical and craft skills (Arcana 3+, Craft 2+).
  ℹtem Costs XP (minimum 5 XP for Tier I) and rare materials.
  ℹtem Takes weeks or months of downtime.
  ℹtem Risk of backlash or permanent curse on failure.
\end{itemize}

For temporary enchantments, use the \textbf{Patron's Gift} (\ref{sec:patrons−gift}) or \textbf{Borrowed Grace} (\ref{sec:invoker−talents}).

\section{Optional Magic Expansions}
łabel{sec:optional−magic}
∈dex{magic!optional systems}

The following systems are optional and may be introduced by the GM.

\begin{tcolorbox}[colback=blue!5!white,colframe=blue!75!black,title=Psionics (Optional)]
∈dex{Psionics}
Psionics represents mental power drawn from within — no Patron, no Symbol, no debt. It is meant to feel \textbf{alien}, even to other magic‑users. Different cultures have varying opinions (see the Grey Wanderer's Grimoire for examples), but psions are generally viewed with suspicion or fear.

\medskip
\textbf{Example Arts:} Telekinesis, Telepathy, Clairvoyance, Biofeedback, Astral Projection, Psychic Assault, Mind Shield, Empathic Manipulation, Precognition.

\medskip
\textbf{Warning:} Psionics generates Story Beats (SB) quickly. The GM is encouraged to spend these SB on \textit{social complications} — suspicion, hostile glances, a Chain‑Lantern investigation, or a neighbor locking their door. Mental backlash is common, but the deepest cost is always \textit{trust}.

\medskip
Each use costs Mental Strain. Overflow converts to normal Fatigue or Harm. Many societies distrust mind‑readers; a psion is always \textit{unwitnessed}, and therefore outside the ledger.
\end{tcolorbox}

\begin{tcolorbox}[colback=green!5!white,colframe=green!75!black,title=Witchcraft (Optional)]
∈dex{Witchcraft}
Witchcraft is magic practiced through patron‑bound Orders. Witches suffer \textbf{—1 die} in consecrated spaces and must navigate the Curse of the Unanointed — the world's refusal to witness their labor.

\medskip
\textbf{Ritual Magic (Weave & Cast):} Two‑step process, each roll generates SB.

\medskip
\textbf{Example Orders:}
\begin{itemize}
  ℹtem \textbf{Mab, Queen of Courts:} Glamour, bargains, fae oaths. Risk — Glamour‑Bloom (lose a memory for +1 Effect).
  ℹtem \textbf{Aveh, The Rider:} Identity shedding. Risk — Unface (your features blur).
  ℹtem \textbf{Morag, The Hag:} Hidden costs and prophecy. Risk — Mirecloth (sacrifice identity for truth).
\end{itemize}

\medskip
\textbf{Identity Cost:} Every use marks who you are. The cost is memory, identity, or reputation.
\end{tcolorbox}

\section{Quick Reference}
łabel{sec:magic−quickref}

\begin{tcolorbox}[colback=black!3!white,colframe=black!40!white,title=Magic Quick Reference,fonttitle=\bfseries]
\textbf{Free Caster}
\begin{itemize}
  ℹtem Talent: Spellcraft (6 XP)
  ℹtem DV = 1 + number of \texttt{[TAGS]}
  ℹtem Backlash on miss or mixed
\end{itemize}

\textbf{Runekeeper}
\begin{itemize}
  ℹtem Thiasos + Codex (6 XP)
  ℹtem 1 action, mark Rite's Obligation cost (variable)
  ℹtem Push It: +1 Obligation
  ℹtem DV = max(Obligation Cost − Spirit, Rite Tier)
\end{itemize}

\textbf{Invoker}
\begin{itemize}
  ℹtem Patron's Symbol (4 XP each)
  ℹtem Ritual: Significant time, mark Rite's Obligation cost
  ℹtem Symbol displayed: reduce Obligation cost by 1
  ℹtem Crack the Seal: instant, double cost (min +2)
\end{itemize}

\textbf{Cantor}
\begin{itemize}
  ℹtem Cantor's Path (8 XP)
  ℹtem Song: 1 action (resolves next round)
  ℹtem Corruption Timer (size = Spirit)
  ℹtem Push: Fatigue +1, advance Corruption
\end{itemize}

\textbf{Summoner}
\begin{itemize}
  ℹtem Pact−Whisperer (2 XP)
  ℹtem Call + Bind (1 Boon or Fatigue)
  ℹtem Leash = Cap + Spirit
\end{itemize}
\end{tcolorbox}